\newpage
\section{Fazit und Ausblick}\label{ausblick}

Zum Abschluss werden die Forschungsfragen aus der Einleitung beantwortet und die Befunde in das Projekt eingeordnet, aus dem die Daten stammen. Daran schließen sich die Ansatzpunkte an, die sich aus den Grenzen des Verfahrens ergeben.
% Zusammenfassung, Limitationen, zukünftige Forschungsmöglichkeiten und
% Einordnung der Ergebnisse in den größeren Kontext.

% --- Muster für eine Abbildung (auskommentiert) ---
% \begin{figure}[H]
% \caption{Einordnung in den Gesamtkontext}\label{fig:ausblick}
% \includegraphics[width=0.9\textwidth]{dateiname}
% \\
% Quelle: Eigene Darstellung
% \end{figure}

\subsection{Fazit}\label{fazit}

Ziel dieser Arbeit war es, den Verbuschungsgrad von Heideflächen objektiv,
flächendeckend und wiederholbar aus dreidimensionalen SfM-Punktwolken
abzuleiten. Dafür wurde PointNet++ auf 16 annotierten Teilflächen trainiert, in
neun Läufen einer Ablationsstudie angepasst und auf vier Befliegungen im
Authausener Wald angewendet.

Auf die erste Forschungsfrage nach der Trennbarkeit der maßgeblichen
Vegetationsklassen gibt es keine einheitliche Antwort. Gehölz und offene
Vegetation trennt das Modell verlässlich, mit einer IoU von 0,761 und einem
F1-Wert von 0,864. Innerhalb der niedrigen Vegetation gelingt die Trennung nicht. Von acht geprüften Eingriffen hat allein die Anpassung der Nachbarschaftsradien an die gemessene Punktdichte gewirkt.

Der automatisierte Label-Transfer der ersten Teilfrage gelingt und macht
den Trainingsdatensatz erst praktikabel, erzeugt jedoch systematische Fehler
an genau den Klassengrenzen, an denen das Modell scheitert.

Ähnlich zweigeteilt fällt die Antwort auf die zweite Forschungsfrage nach der
Übereinstimmung mit der Geländekartierung aus. Auf den sieben Heidepolygonen ohne
Pflegevermerk beträgt der mittlere Fehler 9,5 Prozentpunkte bei einem
Bestimmtheitsmaß von 0,65. Auf der am stärksten verbuschten Fläche stehen
53,0 berechnete gegen 50,0 kartierte Prozent, obwohl das Modell dort nie
trainiert hat. Die Reihenfolge der schwach verbuschten Flächen trifft das
Verfahren dagegen nicht.

Für die zweite Teilfrage bleibt festzuhalten, dass sich die Ergebnisse
unmittelbar in eine bestehende GIS-Umgebung übernehmen lassen, die Inferenz
selbst aber vorerst außerhalb davon läuft.

Damit liegt der in der Einleitung angestrebte durchgängige Weg von der
Befliegung bis zum Verbuschungsgrad vor. Für das Monitoring des
Erhaltungszustands bietet er eine wiederholbare Grundlage, auf der sich
Pflegemaßnahmen planen und ihr Erfolg später überprüfen lassen. Eine
Geländekartierung ersetzt er noch nicht. Wo die verbleibenden Ansatzpunkte liegen, zeigt der folgende Abschnitt.

\subsection{Ausblick und Einordnung}\label{ausblick_teil}

Noch unverarbeitet bleibt das Gebiet der Cuxhavener Küstenheide. Punktwolke und Biotopkartierung liegen vor, doch die Inferenz über 415 Millionen Punkte war im Rahmen dieser Arbeit zeitlich nicht mehr zu leisten. Sie nachzuholen wäre der gewinnbringendste nächste Schritt, denn die Fläche würde die Zahl der Referenzpolygone deutlich erhöhen und zugleich das Ungleichgewicht zwischen Bewertung und Anwendung auflösen.

Der Vergleich mit dem zweidimensionalen Ansatz, der im Projekt parallel entsteht, war ursprünglich Teil dieser Arbeit und musste entfallen, weil dessen Ergebnisse nicht rechtzeitig vorlagen. Nachgeholt werden sollte er auf denselben Kartierungspolygonen und mit denselben Kennzahlen, denn erst dann lässt sich beantworten, ob die dritte Dimension den erheblichen Mehraufwand rechtfertigt.

Methodisch am vielversprechendsten wäre eine objektbezogene statt punktweise Labelkorrektur. Dieses Verfahren würde zusammenhängende Strukturen erkennen und als Ganzes einer Klasse zuweisen, statt sie an einer festen Höhe zu durchschneiden. An ihm ließe sich zugleich die These prüfen, nach der ein Teil der Verwechslung zwischen Heide und Busch aus der Labelvergabe stammt. Daneben würden weitere annotierte Teilflächen den Validierungsteil so weit verbreitern, dass er Modellvarianten wieder unterscheiden kann.

Für den Weg in die Praxis fehlt schließlich ein Gegenstück zu Deepness, das auf Punktwolkenlayern arbeitet. Solange es das nicht gibt, bleibt als naheliegende Lösung, die vorhandenen Skripte als Werkzeug in die Verarbeitungsumgebung von QGIS einzuhängen und die Rechenlast auf einen Server auszulagern, sodass mit dem Ergebnis gearbeitet werden kann, ohne die Verarbeitungskette selbst zu betreiben.

Die Daten dieser Arbeit stammen aus dem in der Einleitung genannten Projekt, das die UAS-gestützte Fernerkundung in das Monitoring von FFH-Gebieten überführen will. Für dessen Fragestellung deckt der hier entwickelte Weg einen von mehreren Degradationsindikatoren ab. Die übrigen, etwa der Deckungsgrad von Gras oder von invasiven Neophyten, ließen sich aus derselben Segmentierung ableiten, sobald das Netz die niedrigen Vegetationsklassen zuverlässiger trennt.
